\section{Heuristiken}

\subsection{Materialistische Evaluation}
Für die materialistische Evaluation wurde die Standardvariation der Basiswerte der Schachfiguren verwendet, jedoch mit 100 multipliziert, um weitere Verfeinerungen leichter berechnen zu können (\cite{evaluationofpieces}). \\
Somit haben die Schachfiguren folgende Basiswerte:
\begin{itemize}
    \item Bauer - 100
    \item Springer - 300
    \item Läufer - 300
    \item Turm - 500
    \item Dame - 900
\end{itemize}

Da der König nicht geschlagen werden kann, aber dennoch ein Angriffsziel bleiben soll, wurde für diese Figur ein hoher Basiswert von 99 definiert, welcher ebenfalls mit 100 multipliziert wurde.

\lstinputlisting[
    firstline=8,
    lastline=14,
    firstnumber=8,
    caption={Basiswerte der Schachfiguren},
    label={lst:figuren-basiswerte}
]{../src/main/java/ch/hslu/cas/msed/blobfish/base/PieceType.java}

\newpage

Für die Berechnung des Materialwertes einer Position werden die Basiswerte aller Figuren aufsummiert.
Die Basiswerte der schwarzen Figuren werden negiert.
Bei einer positiven Endsumme hat der Spieler mit den weissen Figuren den Vorteil, bei einer negativen Summe hat der Spieler mit den schwarzen Figuren den Vorteil. Ist das Ergebnis 0, so ist die Position in Bezug auf das Material ausgeglichen.
\lstinputlisting[
    firstline=15,
    lastline=33,
    firstnumber=15,
    caption={Materialistische Evaluation},
    label={lst:material-evaluation}
]{../src/main/java/ch/hslu/cas/msed/blobfish/eval/MaterialEvaluation.java}

\subsection{Piece Square Tabellen}
Wie bereits erwähnt, bestehen mehrere Wertdefinitionen für Schachfiguren (\cite{evaluationofpieces}).
Genau so existieren auch mehrere \glspl{glos:pst}.
Zunächst sollten die Stockfish \glspl{glos:pst} verwendet werden (\cite{piecesquaretablelichessstockfish}), jedoch wurde festgestellt, dass Stockfish nicht die gleichen Basiswerte der Figuren verwendet (\cite{piecevalueslichessstockfish}).
Die von Stockfish verwendeten Basiswerte der Figuren erschienen unlogisch und hängen wahrscheinlich mit mehreren anderen Faktoren zusammen.

Schlussendlich wurden die \glspl{glos:pst} von \href{https://www.chessprogramming.org/Simplified_Evaluation_Function}{Chess Programming Wiki}~verwendet (\cite{piecesquaretablesef}).
Dementsprechend wurden für den König zwei \glspl{glos:pst} erstellt.
Eine für Eröffnung und das Mittelspiel und eine für das Endspiel.

\newpage

Die Unterscheidung zwischen dem Mittelspiel und dem Endspiel ist ebenfalls nicht standardisiert (\cite{chessendgame}).
Hierfür wurde die gleiche Logik wie bei \gls{glos:lichess} verwendet, welche besagt, dass die Partie sich im Endspiel befindet, sobald nicht mehr als 6 Schachfiguren im Spiel sind.
Die Bauern und Könige sind davon ausgeschlossen:

\lstinputlisting[
    firstline=105,
    lastline=114,
    firstnumber=105,
    caption={Endspiel-Berechnung},
    label={lst:endspiel-berechnung}
]{../src/main/java/ch/hslu/cas/msed/blobfish/base/ChessBoard.java}

Die Implementation der \gls{glos:pst} Evaluation ist oberflächlich gesehen trivial und befindet sich im folgenden Code-Abschnitt.
\lstinputlisting[
    firstline=9,
    lastline=38,
    firstnumber=9,
    caption={Piece Square Evaluation},
    label={lst:piece-square-evaluation}
]{../src/main/java/ch/hslu/cas/msed/blobfish/eval/PieceSquareEvaluation.java}

\newpage

Es wird zuerst ein Array der aktuellen Position generiert.
Die Abbildung~\ref{fig:index-felder-tabelle} stellt das genaue Index-Schachfeld mapping dar.

\begin{figure}[H]
    \centering
    \begin{tabular}{>{\centering\arraybackslash}m{0.2cm}|*{8}{>{\centering\arraybackslash}m{0.4cm}|}}
        \hhline{~|--------}
        8                    & \cellcolor{white} 56   & \cellcolor{gray!30} 57 & \cellcolor{white} 58   & \cellcolor{gray!30} 59 & \cellcolor{white} 60 & \cellcolor{gray!30} 61 & \cellcolor{white} 62 & \cellcolor{gray!30} 63 \\ \hhline{~|--------}
        7                    & \cellcolor{gray!30} 48 & \cellcolor{white} 49   & \cellcolor{gray!30} 50 & \cellcolor{white} 51 & \cellcolor{gray!30} 52 & \cellcolor{white} 53 & \cellcolor{gray!30} 54 & \cellcolor{white} 55\\ \hhline{~|--------}
        6                    & \cellcolor{white} 40   & \cellcolor{gray!30} 41 & \cellcolor{white} 42   & \cellcolor{gray!30} 43 & \cellcolor{white} 44 & \cellcolor{gray!30} 45 & \cellcolor{white} 46 & \cellcolor{gray!30} 47 \\ \hhline{~|--------}
        5                    & \cellcolor{gray!30} 32 & \cellcolor{white} 33   & \cellcolor{gray!30} 34 & \cellcolor{white} 35 & \cellcolor{gray!30} 36 & \cellcolor{white} 37 & \cellcolor{gray!30} 38 & \cellcolor{white} 39 \\ \hhline{~|--------}
        4                    & \cellcolor{white} 24   & \cellcolor{gray!30} 25 & \cellcolor{white} 26   & \cellcolor{gray!30} 27 & \cellcolor{white} 28 & \cellcolor{gray!30} 29 & \cellcolor{white} 30 & \cellcolor{gray!30} 31 \\ \hhline{~|--------}
        3                    & \cellcolor{gray!30} 16 & \cellcolor{white} 17   & \cellcolor{gray!30} 18 & \cellcolor{white} 19 & \cellcolor{gray!30} 20 & \cellcolor{white} 21 & \cellcolor{gray!30} 22 & \cellcolor{white} 23\\ \hhline{~|--------}
        2                    & \cellcolor{white} 8    & \cellcolor{gray!30} 9  & \cellcolor{white} 10   & \cellcolor{gray!30} 11 & \cellcolor{white} 12 & \cellcolor{gray!30} 13 & \cellcolor{white} 14 & \cellcolor{gray!30} 15 \\ \hhline{~|--------}
        1                    & \cellcolor{gray!30} 0  & \cellcolor{white} 1    & \cellcolor{gray!30} 2  & \cellcolor{white} 3 & \cellcolor{gray!30} 4 & \cellcolor{white} 5 & \cellcolor{gray!30} 6 & \cellcolor{white} 7 \\ \hhline{~|--------}
        % Last row with no borders
        \multicolumn{1}{c}{} & \multicolumn{1}{c}{A}  & \multicolumn{1}{c}{B}  & \multicolumn{1}{c}{C} & \multicolumn{1}{c}{D} & \multicolumn{1}{c}{E} & \multicolumn{1}{c}{F} & \multicolumn{1}{c}{G} & \multicolumn{1}{c}{H} \\
    \end{tabular}
    \caption{Indexwerte für die Schachfelder eines Schachbretts}
    \label{fig:index-felder-tabelle}
\end{figure}

Für jedes Feld, wo sich eine Figur befindet, ist im entsprechenden Index ein \texttt{Piece} Objekt enthalten.
Alle leeren Felder enthalten \texttt{null}.

\lstinputlisting[
    firstline=3,
    lastline=3,
    firstnumber=3,
    caption={Piece record},
    label={lst:piece-record}
]{../src/main/java/ch/hslu/cas/msed/blobfish/base/Piece.java}

Nun wird durch das ganze Brett iteriert und überall wo sich eine Figur befindet, wird deren Wert aus der entsprechenden \gls{glos:pst} ausgelesen.
Gleich wie bei den anderen Evaluationen werden Werte für den schwarzen Spieler negiert.
Alle Werte werden am Schluss aufsummiert.

Die Arrays für die \glspl{glos:pst} wurden so strukturiert, dass der Zuschauer die Perspektive seines Schachbretts hat.
Dies sollte das Erstellen und Anpassen der \glspl{glos:pst} für die Entwickler einfach halten.
Für den weissen Spieler ist Index 0 dieses Arrays das Schachfeld \textbf{a8} und für den schwarzen Spieler das Feld \textbf{h1}.

\lstinputlisting[
    firstline=77,
    lastline=86,
    firstnumber=77,
    caption={PST für den Turm},
    label={lst:pst-fuer-turm}
]{../src/main/java/ch/hslu/cas/msed/blobfish/eval/PieceSquareEvaluation.java}

\newpage

Dieser Array-Aufbau entspricht nicht dem Aufbau des vorherigen Schachbrett-Arrays (Abbildung~\ref{fig:index-felder-tabelle}).
Um die richtigen Werte für den weissen Spieler abzurufen, müssen die \glspl{glos:reihe} und für den schwarzen Spieler die \glspl{glos:linie} gespiegelt werden.

\lstinputlisting[
    firstline=9,
    lastline=17,
    firstnumber=9,
    caption={Schachbrett Spiegelung},
    label={lst:schachbrett-spiegelung}
]{../src/main/java/ch/hslu/cas/msed/blobfish/base/BoardTransformationUtil.java}

Die Resultate der Transformationen auf die Indexe vom ursprünglichem Array aus Abbildung~\ref{fig:index-felder-tabelle} sind in der Abbildung~\ref{fig:index-felder-tabelle-spiegelung} ersichtlich.
In dieser Abbildung stellt das Schachbrett den Aufbau der \glspl{glos:pst} dar.
Die Nummern in den Feldern sind die Indexe für den jeweiligen Spieler um den Wert aus einer \gls{glos:pst} für das entsprechende Feld zu holen.

\begin{figure}[H]
    \centering

    \begin{minipage}{0.45\textwidth}
        \centering
        ~~Indexwerte für den weissen Spieler

        \setlength{\tabcolsep}{0pt}
        \renewcommand{\arraystretch}{1}

        \begin{tabular}{>{\centering\arraybackslash}m{0.4cm}|*{8}{>{\centering\arraybackslash}m{0.8cm}|}}
            \hhline{~|--------}
            8                    & \cellcolor{white}    0    & \cellcolor{gray!30}  1   & \cellcolor{white}    2    & \cellcolor{gray!30}  3 & \cellcolor{white}  4& \cellcolor{gray!30} 5& \cellcolor{white} 6& \cellcolor{gray!30} 7\\ \hhline{~|--------}
            7                    & \cellcolor{gray!30}  8  & \cellcolor{white}    9     & \cellcolor{gray!30}   10 & \cellcolor{white}    11 & \cellcolor{gray!30} 12& \cellcolor{white}13 & \cellcolor{gray!30} 14& \cellcolor{white} 15\\ \hhline{~|--------}
            6                    & \cellcolor{white}    16   & \cellcolor{gray!30}  17  & \cellcolor{white}   18   & \cellcolor{gray!30}  19 & \cellcolor{white} 20 & \cellcolor{gray!30} 21 & \cellcolor{white} 22 & \cellcolor{gray!30} 23\\ \hhline{~|--------}
            5                    & \cellcolor{gray!30}  24 & \cellcolor{white}     25    & \cellcolor{gray!30}   26 & \cellcolor{white}     27& \cellcolor{gray!30}28 & \cellcolor{white} 29& \cellcolor{gray!30}30 & \cellcolor{white} 31\\ \hhline{~|--------}
            4                    & \cellcolor{white}    32   & \cellcolor{gray!30} 33 & \cellcolor{white}    34   & \cellcolor{gray!30}  35 & \cellcolor{white}36 & \cellcolor{gray!30} 37& \cellcolor{white}38 & \cellcolor{gray!30}39\\ \hhline{~|--------}
            3                    & \cellcolor{gray!30}  40 & \cellcolor{white}    41    & \cellcolor{gray!30}  42 & \cellcolor{white}    43 & \cellcolor{gray!30} 44 & \cellcolor{white}  45& \cellcolor{gray!30}46 & \cellcolor{white} 47\\ \hhline{~|--------}
            2                    & \cellcolor{white}    48   & \cellcolor{gray!30}  49  & \cellcolor{white}    50   & \cellcolor{gray!30}  51 & \cellcolor{white} 52 & \cellcolor{gray!30} 53 & \cellcolor{white} 54 & \cellcolor{gray!30} 55\\ \hhline{~|--------}
            1                    & \cellcolor{gray!30}  56 & \cellcolor{white}     57    & \cellcolor{gray!30}  58 & \cellcolor{white}    59 & \cellcolor{gray!30} 60& \cellcolor{white} 61& \cellcolor{gray!30} 62& \cellcolor{white} 63\\ \hhline{~|--------}
            % Last row with no borders
            \multicolumn{1}{c}{} & \multicolumn{1}{c}{A} & \multicolumn{1}{c}{B}  & \multicolumn{1}{c}{C} & \multicolumn{1}{c}{D} & \multicolumn{1}{c}{E} & \multicolumn{1}{c}{F} & \multicolumn{1}{c}{G} & \multicolumn{1}{c}{H} \\
        \end{tabular}
    \end{minipage}
    \hfill
    \begin{minipage}{0.45\textwidth}
        \centering
        ~~Indexwerte für den schwarzen Spieler

        \setlength{\tabcolsep}{0pt}
        \renewcommand{\arraystretch}{1}

        \begin{tabular}{>{\centering\arraybackslash}m{0.4cm}|*{8}{>{\centering\arraybackslash}m{0.8cm}|}}
            \hhline{~|--------}
            1                    & \cellcolor{white} 63   & \cellcolor{gray!30} 62 & \cellcolor{white} 61   & \cellcolor{gray!30} 60 & \cellcolor{white} 59 & \cellcolor{gray!30} 58 & \cellcolor{white} 57 & \cellcolor{gray!30} 56\\ \hhline{~|--------}
            2                    & \cellcolor{gray!30} 55 & \cellcolor{white} 54   & \cellcolor{gray!30} 53 & \cellcolor{white} 52 & \cellcolor{gray!30} 51 & \cellcolor{white} 50 & \cellcolor{gray!30} 49 & \cellcolor{white} 48\\ \hhline{~|--------}
            3                    & \cellcolor{white} 47   & \cellcolor{gray!30} 46 & \cellcolor{white} 45   & \cellcolor{gray!30} 44 & \cellcolor{white} 43 & \cellcolor{gray!30} 42 & \cellcolor{white} 41 & \cellcolor{gray!30} 40\\ \hhline{~|--------}
            4                    & \cellcolor{gray!30} 39 & \cellcolor{white} 38   & \cellcolor{gray!30} 37 & \cellcolor{white} 36 & \cellcolor{gray!30} 35 & \cellcolor{white} 34 & \cellcolor{gray!30} 33 & \cellcolor{white} 32\\ \hhline{~|--------}
            5                    & \cellcolor{white} 31   & \cellcolor{gray!30} 30 & \cellcolor{white} 29   & \cellcolor{gray!30} 28 & \cellcolor{white} 27 & \cellcolor{gray!30} 26 & \cellcolor{white} 25 & \cellcolor{gray!30} 24\\ \hhline{~|--------}
            6                    & \cellcolor{gray!30} 23 & \cellcolor{white} 22   & \cellcolor{gray!30} 21 & \cellcolor{white} 20 & \cellcolor{gray!30} 19 & \cellcolor{white} 18 & \cellcolor{gray!30} 17 & \cellcolor{white} 16\\ \hhline{~|--------}
            7                    & \cellcolor{white} 15   & \cellcolor{gray!30} 14 & \cellcolor{white} 13   & \cellcolor{gray!30} 12 & \cellcolor{white} 11 & \cellcolor{gray!30} 10 & \cellcolor{white} 9 & \cellcolor{gray!30} 8\\ \hhline{~|--------}
            8                    & \cellcolor{gray!30} 7  & \cellcolor{white} 6    & \cellcolor{gray!30} 5  & \cellcolor{white} 4 & \cellcolor{gray!30} 3 & \cellcolor{white} 2 & \cellcolor{gray!30} 1 & \cellcolor{white} 0\\ \hhline{~|--------}
            % Last row with no borders
            \multicolumn{1}{c}{} & \multicolumn{1}{c}{H}  & \multicolumn{1}{c}{G}  & \multicolumn{1}{c}{F} & \multicolumn{1}{c}{E} & \multicolumn{1}{c}{D} & \multicolumn{1}{c}{C} & \multicolumn{1}{c}{B} & \multicolumn{1}{c}{A} \\
        \end{tabular}

    \end{minipage}
    \caption{Indexwerte-Spiegelung für die Schachfelder}
    \label{fig:index-felder-tabelle-spiegelung}
\end{figure}

\textbf{Ein konkretes Beispiel:} \\
Ein weisser Turm befindet sich auf dem Schachfeld \textbf{b2}.
Gemäss Abbildung~\ref{fig:index-felder-tabelle} beträgt der ursprüngliche Index für dieses \gls{glos:feld} 9.
Der benötigte Index um den korrekten Wert aus der \gls{glos:pst} auszulesen ist 49.
Die entsprechende \gls{glos:pst} liefert gemäss Code-Ausschnitt~\ref{lst:pst-fuer-turm} den Wert 0.
Ein schwarzer Turm auf dem gleichen Schachfeld liefert den Wert 10.